\subsection{Mode powers and the Mumford subsystem}
\label{sec:mumford-correspondence}

Throughout this subsection $N=1$, with the hyperelliptic data of
\cref{sec:BA-realization}.

The elementary mode powers in \cref{cor:BA-realization} are
evaluations of the moving divisor polynomial $\mathcal P$. Together
with their derivatives, they determine a closed projected subsystem
in divisor coordinates, which we identify with the odd Mumford flow.
We call this the Mumford subsystem. We reconstruct its polynomials
from the modes and describe the mode information retained by the
BA formulas. The curve, marked points and weights remain fixed throughout this subsection. We write $\mathcal P_{z_g}=\partial_{z_g}\mathcal P=i\dot{\mathcal P}$ and
$\mathcal P_{z_gz_g}=\partial_{z_g}^2\mathcal P=-\ddot{\mathcal P}$.

\subsubsection{Reconstructing the divisor polynomials}

Define the elementary powers $I_k=\xi_k\bar{\xi}_k$, $1\leq k\leq K$.
Their derivatives follow directly from the canonical Hamiltonian
\cref{eq:BA-H0-flat}: the linear phase terms cancel, leaving
\begin{equation}\label{eq:mumford-power-rates}
\begin{aligned}
 i\dot I_{\ell}&=\sqrt2\sum_{j=L+1}^{K}
 (\bar{\xi}_{\ell}\bar{\xi}_j\bar{\xi}_{\ell j}-\xi_{\ell}\xi_j\xi_{\ell j}),&&1\leq\ell\leq L,\\
 i\dot I_j&=\sqrt2\sum_{\ell=1}^{L}
 (\bar{\xi}_{\ell}\bar{\xi}_j\bar{\xi}_{\ell j}-\xi_{\ell}\xi_j\xi_{\ell j}),&&L+1\leq j\leq K.
\end{aligned}
\end{equation}
Thus $\dot I_k=\{I_k,H_0\}$ is a polynomial in the modes, rather
than additional derivative data. On the BA family,
\begin{equation}\label{eq:mumford-evaluations}
 \frac{I_k}{\kappa_k}=\mathcal P(z,a_k),\qquad
 \frac{i\dot I_k}{\kappa_k}=\mathcal P_{z_g}(z,a_k),
 \qquad1\leq k\leq K.
\end{equation}

Since $L\geq g$, use the first $g$ marked abscissae and their
Lagrange basis polynomials
\begin{equation}\label{eq:mumford-lagrange}
 \ell_k(x)=\prod_{\substack{r=1\\r\ne k}}^{g}
                  \frac{x-a_r}{a_k-a_r},\qquad1\leq k\leq g.
\end{equation}
The leading coefficient of $\mathcal P$ is one, so interpolation
of its remaining coefficients and of $\mathcal P_{z_g}$ gives
\begin{equation}\label{eq:mumford-reconstruction}
\begin{aligned}
 \mathcal P(z,x)
 &=x^g+\sum_{k=1}^{g}
       \left(\frac{I_k}{\kappa_k}-a_k^g\right)\ell_k(x),\\
 \mathcal P_{z_g}(z,x)
 &=\sum_{k=1}^{g}\frac{i\dot I_k}{\kappa_k}\ell_k(x).
\end{aligned}
\end{equation}
Any $g$ distinct marked abscissae give the same reconstruction on
this family. The additional evaluations supply the remaining mode
powers and their linear interpolation relations.

Jacobi inversion gives $\mathcal P\mid f-\mathcal P_{z_g}^{\,2}$:
at each moving divisor point $(x_r,y_r)$ one has
$\mathcal P_{z_g}(z,x_r)=y_r$ and hence $f(x_r)=\mathcal P_{z_g}(z,x_r)^2$.
Consequently define the polynomial
\begin{equation}\label{eq:mumford-W}
 \mathcal W(z,x)=\frac{f(x)-\mathcal P_{z_g}(z,x)^2}
                         {4\mathcal P(z,x)},\qquad
 f=\mathcal P_{z_g}^{\,2}+4\mathcal P\mathcal W.
\end{equation}
Here $\mathcal P$ and $\mathcal W$ are monic of degrees $g$ and
$g+1$, while $\deg_x\mathcal P_{z_g}\leq g-1$. The factor four
makes $\mathcal W$ monic for the curve normalization
\cref{eq:BA-curve}. Write Mumford's three polynomials as
$(\mathcal U,\mathcal V,\mathcal W)$, with
\begin{equation}\label{eq:mumford-dictionary}
 \mathcal U=\mathcal P,\qquad
 \mathcal V=\mathcal P_{z_g}/2,\qquad
 \mathcal U\mathcal W+\mathcal V^2=f/4.
\end{equation}
These are the degree constraints and spectral relation of the
ordinary odd Mumford system
\cite[Chapter~IIIa]{Mumford1984}; see also
\cite[\S3.1]{FernandesVanhaecke2001}. Its monic spectral curve is
obtained from ours by dividing the ordinate by two.

\subsubsection{The projected equations of motion}

Write the common scalar potential as
\begin{equation}\label{eq:mumford-potential}
 \mathcal Q(z)=2\wp_{gg}(z)+\frac{\lambda_{2g}}4
 =\frac{\lambda_{2g}}4-2[x^{g-1}]\mathcal P(z,x).
\end{equation}
The shared second-order BA equation and constant sheet-pair
Wronskian in \cref{eq:BA-identities,eq:BA-powers-rates} give
\begin{equation}\label{eq:mumford-product-identity}
 2\mathcal P\mathcal P_{z_gz_g}-\mathcal P_{z_g}^{\,2}
       -4(x+\mathcal Q)\mathcal P^2=-f(x).
\end{equation}
For example, differentiate the product of the BA functions at an
arbitrary curve point with abscissa $x$, use their common scalar
equation, and eliminate the difference of their derivatives with
their Wronskian. Since the point is arbitrary, this gives a
polynomial identity in $x$.

Substituting \cref{eq:mumford-W} into
\cref{eq:mumford-product-identity}, then differentiating the
spectral relation, yields the closed polynomial flow
\begin{equation}\label{eq:mumford-flow}
 \mathcal P_{z_gz_g}=2\bigl((x+\mathcal Q)\mathcal P-\mathcal W\bigr),
 \qquad
 i\dot{\mathcal W}=-(x+\mathcal Q)\mathcal P_{z_g}.
\end{equation}
Together with $\mathcal P_{z_g}$ as the second polynomial
coordinate, these are the odd Mumford infinity flow. The time
normalization is fixed by $i\dot{\mathcal U}=2\mathcal V$
and the BA direction $z(t)=z_0-it e_g$.

The correspondence can also be checked in a two-by-two Lax form.
Define
\begin{equation}\label{eq:mumford-lax-matrices}
 \mathsf A(x)=\begin{pmatrix}
 \mathcal V&\mathcal U\\
 \mathcal W&-\mathcal V
 \end{pmatrix},\qquad
 \mathsf B(x)=\begin{pmatrix}0&1\\x+\mathcal Q&0\end{pmatrix}.
\end{equation}
Then \cref{eq:mumford-flow} is equivalent to
\begin{equation}\label{eq:mumford-lax}
 i\dot{\mathsf A}=[\mathsf A,\mathsf B],\qquad
 \det\left(\frac y2\Id_2-\mathsf A(x)\right)
                  =\frac{y^2-f(x)}4.
\end{equation}
Thus the mode-power projection preserves both the divisor dynamics
and the prescribed hyperelliptic curve.

\subsubsection{Recovering the modes from divisor data}

At a generic state of the fixed-curve Mumford fiber, set
$\mathcal P_k=\mathcal P(a_k)$ and
$\mathcal P_{z_g,k}=\mathcal P_{z_g}(a_k)$,
$1\leq k\leq K$. Choose nonzero elementary values $\xi_k$ and put
\begin{equation}\label{eq:mumford-mode-inverse}
 \bar{\xi}_k=\frac{\kappa_k\mathcal P_k}{\xi_k},
 \qquad1\leq k\leq K.
\end{equation}
Substitution of the BA logarithmic rates into the Wronskians
\cref{eq:BA-pair-composites} determines every composite algebraically:
\begin{equation}\label{eq:mumford-composite-inverse}
\begin{aligned}
 \xi_{\ell j}
 &=\frac{\bar{\xi}_{\ell}\bar{\xi}_j}{\sqrt2(a_j-a_{\ell})}
       \left(\frac{\mathcal P_{z_g,j}+y_j}{\mathcal P_j}
             -\frac{\mathcal P_{z_g,i}+y_{\ell}}{\mathcal P_{\ell}}\right),\\
 \bar{\xi}_{\ell j}
 &=-\frac{\xi_{\ell}\xi_j}{\sqrt2(a_j-a_{\ell})}
       \left(\frac{\mathcal P_{z_g,j}-y_j}{\mathcal P_j}
             -\frac{\mathcal P_{z_g,i}-y_{\ell}}{\mathcal P_{\ell}}\right),
 \qquad1\leq\ell\leq L<j\leq K.
\end{aligned}
\end{equation}
These formulas apply on the chart $\mathcal P_k\ne0$,
$1\leq k\leq K$. Jacobi inversion supplies the corresponding
$z$, and the constants $c_k$ in \cref{eq:BA-elementary} match the
chosen elementary values. Under the conditions of
\cref{cor:BA-realization}, the lifted modes then follow the
canonical mode equations.

The polynomial triple determines the divisor state; the $K$
elementary normalizations specify its mode lift. In particular,
the phase-dependent individual BA functions are retained by the
mode description, and the composites are fixed once those
normalizations have been chosen. After forgetting this gauge
freedom, the fixed-curve BA family gives an open part of the usual
Mumford fiber, a Jacobian with its theta divisor removed
\cite[Proposition~3.1]{FernandesVanhaecke2001}. This identifies
the fixed-curve solutions and their vector fields. A correspondence
of Poisson structures would additionally require a reduction of
the canonical mode bracket.

\subsubsection{Mode gauges and the divisor lift}
\label{sec:mumford-gauge-lift}

We specialize the general symmetry interpretation of
\cref{sec:gauge-Abel-symmetries} to identify the mode phases omitted
by the divisor projection.

The $K=L+M$ independent one-point charges generate the reciprocal
mode gauges
\begin{equation}\label{eq:BA-gauge-action}
\begin{aligned}
 \xi_k&\longmapsto e^{\theta_k}\xi_k,&
 \bar{\xi}_k&\longmapsto e^{-\theta_k}\bar{\xi}_k,
 &&1\leq k\leq K,\\
 \xi_{\ell j}&\longmapsto e^{-(\theta_{\ell}+\theta_j)}\xi_{\ell j},&
 \bar{\xi}_{\ell j}&\longmapsto e^{\theta_{\ell}+\theta_j}\bar{\xi}_{\ell j},
 &&1\leq\ell\leq L<j\leq K.
\end{aligned}
\end{equation}
Here the $\theta_k$ are constant complex parameters; purely imaginary
values preserve the conjugate real form. Every term of $H_0$ is
invariant. These symmetries follow from the polynomial Hamiltonian
itself; the BA realization identifies them with changes of the
normalization constants $c_k$.

For the prescribed hyperelliptic BA family, $g$ divisor coordinates and
$K$ mode normalizations specify the lift described in
\cref{eq:mumford-mode-inverse,eq:mumford-composite-inverse}.
Its spectral level is identified in
\cref{prop:BA-prescribed-spectral-curve}. The realized $H_0$
flow translates $z$ in the $-ie_g$ direction.
\Cref{thm:BA-higher-flows} identifies the further spectral flows
with translations of this same BA family and the constant mode phases
in \cref{eq:BA-higher-time-phase}. Their Abel directions span its
Jacobian, while the power charges generate the reciprocal gauges.

\subsubsection{Mechanical coordinate realizations}

An explicit connection with Garnier coordinates follows by taking
$L=g$. Then $q_L=1$ and $\kappa_{\ell}=1/(2R_L'(a_{\ell}))$,
$1\leq\ell\leq g$. Introduce local coordinates $r_{\ell}$ by
$r_{\ell}^2=4I_{\ell}$. Lagrange interpolation gives
\begin{equation}\label{eq:mumford-garnier-potential}
 \sum_{\ell=1}^{g}r_{\ell}^2
 =2\left([x^{g-1}]\mathcal P+\sum_{\ell=1}^{g}a_{\ell}\right).
\end{equation}
Using \cref{eq:mumford-product-identity} to differentiate these
coordinates gives
\begin{equation}\label{eq:mumford-garnier}
 -\ddot r_{\ell}
 =\left(\alpha_{\ell}-\sum_{j=1}^{g}r_j^2\right)r_{\ell}
   -\frac{f(a_{\ell})}{R_L'(a_{\ell})^2r_{\ell}^3},\qquad
 \alpha_{\ell}=a_{\ell}+\frac{\lambda_{2g}}4+2\sum_{j=1}^{g}a_j,
 \quad1\leq\ell\leq g.
\end{equation}
In the geometric Abel time $-it$, this is the Garnier form with centrifugal terms
\cite{Kostov1999,Vanhaecke2001}. At marked branch points the
corresponding inverse-cube terms vanish. The elementary BA modes
provide a description retaining the sheet-pair phases as well as
these mechanical coordinates.

For the Neumann realization of the same divisor geometry, see
\cite[Chapter~IIIa, \S4]{Mumford1984}. Vanhaecke discusses its
connections with Garnier and integrable H\'enon--Heiles systems
\cite{Vanhaecke2001}; the genus-two case~(ii) H\'enon--Heiles
system and its BA description are treated explicitly in
\cite[\S4]{Braden2005}.
